% ==============================================================
%  INSTRUMENTO 4 — FICHA DE OBSERVACIÓN
% ==============================================================

\textbf{Resultados del instrumento 4}

A partir del diagnóstico cualitativo obtenido en las entrevistas semiestructuradas (Instrumento~1), se determinó que la actualización y transcripción de historias clínicas no se ejecuta durante la atención presencial del paciente, sino de forma diferida en bloque al finalizar la jornada o en momentos de baja afluencia. Por consiguiente, el Instrumento~4 se estructuró en dos formatos independientes (Anexo~2): la Parte~A (medición en tiempo real por paciente) y la Parte~B (registro de lote diario de digitalización).

\textbf{Resultados de la Parte~A: Medición por paciente ($n = 87$)}

La medición directa se aplicó a 87 pacientes en jornadas de alta afluencia (Martes, Jueves y Sábado). El tiempo administrativo directo por paciente alcanzó un promedio de 9{,}4~minutos. Adicionalmente, se registraron tiempos de espera previa y retraso respecto a la cita programada derivados de la gestión manual. Los resultados se resumen en la Tabla~\ref{tab:resumen_obs_parte_a}.

\begin{table}[H]
    \centering
    \fontsize{10}{12}\selectfont
    \caption{Promedios de tiempos administrativos por paciente en fase pretest (Parte~A, $n = 87$)}
    \label{tab:resumen_obs_parte_a}
    \renewcommand{\arraystretch}{1.2}
    \begin{tabularx}{\linewidth}{|l|X|c|}
        \hline
        \textbf{Tramo / Indicador} & \textbf{Descripción} & \textbf{Promedio (min)} \\ \hline
        $T_4$ — Registro post-consulta & Elaboración manual de receta y anotación del procedimiento & 4,5 \\ \hline
        $T_5$ — Registro de pago & Registro del cobro, abono y emisión de comprobante & 2,2 \\ \hline
        $T_6$ — Agendamiento próxima cita & Coordinación y asignación de turno en Google Calendar & 2,7 \\ \hline
        \textbf{Tiempo admin. directo} & \textbf{Suma de tramos ($T_4 + T_5 + T_6$)} & \textbf{9,4} \\ \hline
        Espera previa a la atención & Tiempo transcurrido entre la llegada y el inicio en sillón & 11,4 \\ \hline
        Retraso vs. cita programada & Desfase entre la hora agendada y el inicio de la atención & 8,4 \\ \hline
    \end{tabularx}
\end{table}

Entre las incidencias observadas con mayor frecuencia destacaron la redacción manual de prescripciones médicas con consulta física de códigos CIE-10, la búsqueda tardía de expedientes en pacientes sin cita previa y desajustes en el agendamiento por confirmaciones informales vía WhatsApp.

\textbf{Resultados de la Parte~B: Registro de lote diario de digitalización}

La digitalización diferida se evaluó durante tres jornadas completas, registrando la transcripción de 86 historias clínicas físicas (17 pacientes nuevos y 69 recurrentes). La tarea demandó un acumulado de 330~minutos (110~minutos por día). El tiempo promedio de digitalización unitaria fue de 3{,}84~minutos por paciente, tal como se detalla en la Tabla~\ref{tab:resumen_obs_parte_b}.

\begin{table}[H]
    \centering
    \fontsize{10}{12}\selectfont
    \caption{Registro de lote diario para digitalización de historias clínicas en fase pretest (Parte~B)}
    \label{tab:resumen_obs_parte_b}
    \renewcommand{\arraystretch}{1.2}
    \begin{tabularx}{\linewidth}{|X|c|c|c|c|}
        \hline
        \textbf{Jornada} & \textbf{Fecha} & \textbf{Duración (min)} & \textbf{Pacientes} & \textbf{Promedio / Paciente (min)} \\ \hline
        Martes (Vespertino) & 22/09/2026 & 110 & 29 & 3,79 \\ \hline
        Jueves (Matutino) & 24/09/2026 & 105 & 28 & 3,75 \\ \hline
        Sábado (Vespertino) & 26/09/2026 & 115 & 29 & 3,97 \\ \hline
        \textbf{Total / Promedio} & \textbf{3 días} & \textbf{330} & \textbf{86} & \textbf{3,84} \\ \hline
    \end{tabularx}
\end{table}

Durante la digitalización se observaron interrupciones frecuentes (hasta 15~minutos por jornada) causadas por la atención simultánea de llamadas telefónicas y mensajes en recepción, así como demoras puntuales al descifrar grafías manuscritas en las fichas físicas.


